\documentclass[10pt]{beamer}
\usetheme{Madrid}
\usepackage{booktabs}
\usepackage{graphicx}
\usepackage[T1]{fontenc}
\setbeamertemplate{navigation symbols}{}

% ======================= EDIT YOUR DETAILS =======================
\newcommand{\studentname}{Aflah Muhammed P}
% Change the university roll/registration number:
\newcommand{\rollno}{TVE23CSXXX}
% Change your class roll number:
\newcommand{\rollclass}{Roll No: XX}
% Change your guide name:
\newcommand{\guidename}{Prof. Guide Name}
% Change department / college / short-institute if needed:
\newcommand{\deptname}{Department of Computer Science and Engineering}
\newcommand{\collegename}{College of Engineering Trivandrum}
\newcommand{\shortinst}{CET}
% Change the seminar date:
\newcommand{\seminardate}{\today}
% =================================================================

\title[B.Tech Seminar]{How AI Agents Think: A Beginner's Map of LLM Agentic Reasoning Frameworks}
\subtitle{Based on: LLM-based Agentic Reasoning Frameworks: A Survey from Methods to Scenarios (arXiv:2508.17692, 2025)}
\author[\studentname\ (\shortinst)]{\studentname \\ \small \rollno \\ \small \rollclass \\ \small Guide: \guidename}
\institute[]{\deptname \\ \collegename}
\date{\seminardate}

\begin{document}

\begin{frame}[plain]
  \titlepage
\end{frame}

\begin{frame}{Table of Contents}
  \tableofcontents
\end{frame}

\section{Introduction}
\begin{frame}{Introduction}
\begin{itemize}
  \item LLMs alone give one answer; agents act in a loop
  \item AI agents now power coding, healthcare, and science tools
  \item Many competing reasoning designs, but no shared vocabulary
  \item Survey builds one map from methods to real scenarios
  \item Goal: understand each design's strengths, fit, and evaluation
\end{itemize}
\end{frame}

\section{Literature Survey}
\begin{frame}{Literature Survey / Background}
  \small
  \begin{tabular}{@{}p{2.7cm}p{2.3cm}p{5.6cm}@{}}
    \toprule
    \textbf{Title} & \textbf{Author} & \textbf{Summary} \\
    \midrule
    ReAct: Reasoning and Acting in Language Models & Shunyu Yao et al. & Combines step-by-step reasoning with actions so an agent thinks, calls a tool, observes the result, then continues toward its goal. \\
    \addlinespace
    Chain-of-Thought Prompting & Jason Wei et al. & Prompts a model to show its intermediate reasoning steps, sharply improving performance on math and logic problems. \\
    \addlinespace
    MetaGPT: Multi-Agent Collaborative Framework & Sirui Hong et al. & Assigns software-team roles to multiple agents so they collaborate to turn one requirement into working code. \\
    \addlinespace
    \bottomrule
  \end{tabular}
\end{frame}

\section{Motivation}
\begin{frame}{Motivation}
\begin{itemize}
  \item Agent research scattered across hundreds of inconsistent papers
  \item Hard to compare frameworks without a shared language
  \item Beginners cannot tell which design fits which task
  \item Evaluation practices differ widely across domains
\end{itemize}
\end{frame}

\section{Problem Statement}
\begin{frame}{Problem Statement}
  \begin{block}{}
  There is no unified way to describe, categorize, and compare the many LLM-based agentic reasoning frameworks, making it hard to know which design suits a given task and how to evaluate it fairly. This survey addresses that gap with a common formal language and a clear taxonomy.
  \end{block}
\end{frame}

\section{Objectives}
\begin{frame}{Objectives}
\begin{itemize}
  \item Introduce a unified formal language for describing agents
  \item Build a taxonomy of agentic reasoning frameworks
  \item Classify methods as single-agent, tool-based, or multi-agent
  \item Map each method family to real application scenarios
  \item Summarize evaluation strategies across the domains
\end{itemize}
\end{frame}

\section{Methodology}
\begin{frame}[allowframebreaks]{Methodology / Approach}
\begin{itemize}
  \item Model an agent as a context updated by actions
  \item Define an action space: reason, tool-use, and reflect
  \item Loop until the goal or a stopping condition is met
  \item Group all frameworks into three main families
  \item Break each family into concrete sub-types
  \item Connect the families to five real-world domains
\end{itemize}
\end{frame}

\section{Key Concepts}
\begin{frame}[allowframebreaks]{Key Concepts}
  \begin{description}
    \item[Single-agent methods] One agent reasons alone, using prompt-engineering tricks and self-improvement like reflection and iterative optimization.
    \item[Tool-based methods] Agents call external tools; covers how tools are integrated, selected, and used in sequence or parallel.
    \item[Multi-agent methods] Multiple agents organized centrally or hierarchically, interacting through cooperation, competition, or negotiation.
    \item[Unified formal language] A simple shared notation describing any agent as a context plus reason, tool, and reflect actions.
    \item[Application scenarios] Five domains studied: science, healthcare, software engineering, social simulation, and economics.
  \end{description}
\end{frame}

\section{System Architecture}
\begin{frame}{System Architecture / How It Works}
  \begin{block}{Overview}
  The core pipeline is a loop. An agent starts with a context that holds the user goal, then repeatedly chooses an action of one of three types: reason (think), tool-use (call an external resource), or reflect (critique itself). Each action updates the context, and the loop continues until a termination condition is satisfied and the answer is returned. A diagram can show one central Context box with arrows cycling out to Reason, Tool, and Reflect blocks and back, then branching into the three framework families (single-agent, tool-based, multi-agent) and finally into the five application domains.
  \end{block}
  % TIP: insert a diagram here, e.g.:
  % \begin{center}\includegraphics[width=0.7\textwidth]{your-figure.png}\end{center}
\end{frame}

\section{Results}
\begin{frame}[allowframebreaks]{Results / Findings}
\begin{itemize}
  \item Three-family taxonomy unifies scattered agent designs
  \item No single framework is best for every task
  \item Tool use expands ability but adds selection complexity
  \item Multi-agent suits collaborative tasks at higher cost
  \item Survey spans 51 pages, 10 figures, and 8 tables
\end{itemize}
\end{frame}

\section{Discussion}
\begin{frame}{Discussion}
\begin{itemize}
  \item Shared language makes different frameworks directly comparable
  \item Choice of family should match task complexity and cost
  \item Explicit tool and action steps clarify safety and failure points
  \item Bridges method-level research and real deployment needs
\end{itemize}
\end{frame}

\section{Challenges}
\begin{frame}{Limitations \& Challenges}
\begin{itemize}
  \item A survey, not new experiments or head-to-head benchmarks
  \item Fast-moving field; new frameworks appear constantly
  \item Evaluation practices remain inconsistent across domains
  \item Real-world reliability and safety are still open problems
\end{itemize}
\end{frame}

\section{Future Work}
\begin{frame}{Future Work}
\begin{itemize}
  \item Standardize evaluation across agent frameworks
  \item Improve reliability, cost, and coordination of agents
  \item Extend the taxonomy as new methods emerge
  \item Strengthen safety and robustness in deployment
\end{itemize}
\end{frame}

\section{Conclusion}
\begin{frame}{Conclusion}
\begin{itemize}
  \item Provides a clear map of LLM agent reasoning
  \item Three families: single-agent, tool-based, and multi-agent
  \item A shared formal language enables fair comparison
  \item Links methods to real scenarios and their evaluation
\end{itemize}
\end{frame}

\section{References}
\begin{frame}[allowframebreaks]{References}
  \footnotesize
  \begin{enumerate}
    \item LLM-based Agentic Reasoning Frameworks: A Survey from Methods to Scenarios, Bingxi Zhao, Lin Geng Foo, Ping Hu, Christian Theobalt, Hossein Rahmani, Jun Liu, arXiv:2508.17692, 2025
    \item ReAct: Synergizing Reasoning and Acting in Language Models, Shunyu Yao, Jeffrey Zhao, Dian Yu, Nan Du, Izhak Shafran, Karthik Narasimhan, Yuan Cao, ICLR, 2023
    \item Chain-of-Thought Prompting Elicits Reasoning in Large Language Models, Jason Wei et al., NeurIPS, 2022
    \item Reflexion: Language Agents with Verbal Reinforcement Learning, Noah Shinn et al., NeurIPS, 2023
    \item MetaGPT: Meta Programming for a Multi-Agent Collaborative Framework, Sirui Hong et al., ICLR, 2024
    \item Tree of Thoughts: Deliberate Problem Solving with Large Language Models, Shunyu Yao et al., NeurIPS, 2023
  \end{enumerate}
\end{frame}

\begin{frame}[plain]
  \begin{center}
    {\Huge Thank You}\\[1em]
    {\large Questions?}
  \end{center}
\end{frame}

\end{document}